% !TEX root = ../main.tex
\section{Conclusiones}

Implementamos y evaluamos el Splay Tree y el bosque dinámico Link-Cut Tree en Python, verificando experimentalmente su comportamiento amortizado, su corrección estructural y su rendimiento en diferentes tipos de grafos.

En los ejercicios resueltos comprobamos que las rotaciones locales conservan el orden simétrico en árboles de búsqueda binaria y aplanan configuraciones lineales en tiempo constante por paso. En el Link-Cut Tree diferenciamos con claridad el árbol representado respecto a los árboles auxiliares, confirmando que las rotaciones dentro de los caminos preferidos nunca alteran las aristas lógicas del bosque. El uso de la marca diferida de inversión permitió re-enraizar componentes y exponer caminos manteniendo sumas acumuladas y tamaños de subárbol correctos.

La validación diferencial contrastó 15\,000 operaciones aleatorias frente a un modelo de referencia basado en búsqueda en anchura, alcanzando cero discrepancias y verificando nueve invariantes estructurales en cada paso. El sistema interceptó correctamente los intentos de crear ciclos y los cortes sobre aristas inexistentes o no directas.

En las pruebas de rendimiento con tamaños de hasta 1\,600 nodos, la cadena lineal generó el mayor número de rotaciones por operación debido a la reestructuración de caminos profundos, mientras que la estrella mantuvo un costo constante. Concluimos que, aunque una operación aislada puede costar tiempo lineal en el peor caso, la liberación de potencial acumulado asegura un costo amortizado logarítmico para cualquier secuencia de operaciones. Finalmente, observamos que el costo de manejar objetos y punteros en Python penaliza el tiempo de ejecución real frente a bibliotecas nativas, demostrando que la cota asintótica no equivale directamente a la velocidad observada.
